\section{Satisfiability}
\label{sec:ch09-satisfiability}
\index{satisfiability|(}%

Two schemas can merge cleanly into a schema that describes a graph nobody can
query. That is the failure satisfiability exists to catch, and it is the one
composition check with no equivalent in a single-service world.

Apollo states it as one of three composition rules: \enquote{If multiple
subgraphs define the same type, each field of that type must be resolvable by
every valid GraphQL operation that includes it}, and calls it \enquote{the most
complex and the most essential to Federation
2}~\autocite{apollo2026compositionrules}.

Read that twice, because the quantifier is the whole rule. Not \emph{some}
operation. Every valid operation that includes the field. A graph is
satisfiable when there is no query a client could write that the router would
have to give up on.

\subsection*{It is a walk}

\index{satisfiability!as a graph walk}%
Cosmo implements the rule as a graph traversal, in a directory called
\code{resolvability-graph}. The names in it are the mechanism:
\code{graph-nodes}, \code{node-resolution-data}, \code{walker}.

Figure~\ref{fig:ch09-resolvability-walk} is the shape of it for Mosaic. Start at
a root field. That field belongs to one subgraph, so you are now inside that
subgraph and can reach anything it declares. To reach a field it does not
declare, you have to leave, and the only way out is a key on an entity you are
currently holding, pointing at a subgraph that will accept it. Every field the
walk never arrives at is a satisfiability error.

\begin{figure}[htbp]
  \centering
  % Satisfiability as the walk it is: from a root field, through the subgraph
% that owns it, across a key edge, to the fields another subgraph contributes.
% Take the key edge away and the four fields on the right are unreachable,
% which is what chapter 09 measures. Referenced from chapter 09. Bare
% tikzpicture; the figure environment, caption and label live at the call site.
%
% The zone divider is drawn only above the key edge. Run it the full height and
% it crosses the one line in the picture that matters.
\begin{tikzpicture}[
    font=\small,
    place/.style={draw, rounded corners=2pt, minimum width=34mm,
                  minimum height=12mm, align=center, font=\scriptsize},
    entry/.style={draw, rounded corners=2pt, minimum width=34mm,
                  minimum height=8mm, align=center, font=\scriptsize,
                  fill=black!8},
    walk/.style={-{Stealth[length=2mm]}},
    keyedge/.style={-{Stealth[length=2mm]}, thick},
    cut/.style={thick},
    lbl/.style={font=\scriptsize, align=center, inner sep=1pt},
    onedge/.style={font=\scriptsize, align=center, fill=white, inner sep=2pt},
    zone/.style={gray, dashed, thick}
  ]

  % the two subgraphs, named above the boxes rather than boxed themselves
  \draw[zone] (6.5,1.6) -- (6.5,0.05);
  \node[lbl, text=gray, anchor=south] at (2.9,1.65) {subgraph \code{catalog}};
  \node[lbl, text=gray, anchor=south] at (10.1,1.65) {subgraph \code{mosaic}};

  \node[entry] (root) at (2.9,1.0) {\code{Query.products}};
  \node[place] (pcat) at (2.9,-0.9)
    {\code{Product} here\\\code{id sku title}\\\code{description category}};

  \node[place] (pmos) at (10.1,-0.9)
    {\code{Product} here\\\code{price availableQuantity}\\\code{reviews averageRating}};

  \draw[walk] (root) -- (pcat);
  \draw[keyedge] (pcat) -- (pmos);
  \node[onedge] at (6.5,-0.45) {\code{@key(fields: "id")}};

  % what resolvable: false does to it
  \draw[cut] (6.28,-0.68) -- (6.72,-1.12);
  \draw[cut] (6.28,-1.12) -- (6.72,-0.68);
  \node[lbl, anchor=north, align=center] at (6.5,-1.35)
    {\code{resolvable: false}\\removes this edge};

\end{tikzpicture}

  \caption{The only way out of a subgraph is a key another subgraph will answer
  for. Every field the walk never reaches is a field some client query can ask
  for and no router can answer, which is what the composer refuses to emit. One
  edge carries all four of Mosaic's fields here, so cutting it produces four
  errors rather than one.}
  \label{fig:ch09-resolvability-walk}
\end{figure}

\index{satisfiability!error factories}%
There is exactly one error factory for all of this. In
\code{@wundergraph/composition} at 0.63.2, \code{unresolvablePathError} is the
only one that satisfiability produces:

\begin{minted}[fontsize=\footnotesize]{javascript}
function unresolvablePathError({ fieldName, selectionSet }, reasons) {
    const message = `The field "${fieldName}" is unresolvable at the following path:\n${selectionSet}` +
        `\nThis is because:\n - ` +
        reasons.join(`\n - `);
    return new Error(message);
}
\end{minted}

Everything the check can tell you is in that \code{reasons} array, which two
functions in the same package assemble. Reading those two is the cheapest way to
learn what the walk knows about itself: which root field it started from, which
subgraphs define the field, which of them marked it \code{@external}, and
whether the entity it was standing on could satisfy a key to reach the subgraph
that has what it wanted.

\subsection*{Breaking it on purpose}

\index{satisfiability!measured on Mosaic}%
The way to remove that edge without removing anything else is one argument.
\code{@key} takes a \code{resolvable} flag, and a key declared
\code{resolvable: false} means the type is still an entity here, still carries
the key, and still may not be entered with it.

So Mosaic's \code{Product} becomes:

\begin{graphqlsdl}
type Product @key(fields: "id", resolvable: false) {
  availableQuantity: Int!
  price: Money!
  reviews(first: Int, after: String, last: Int, before: String): ReviewConnection!
  averageRating: Float
  id: ID!
}
\end{graphqlsdl}

One argument, nothing else touched. The schema is still valid, Mosaic would
still start, and the merge still produces a sensible-looking \code{Product} with
nine fields on it. The composer refuses, four times, once per field Mosaic
contributes. Here is the first, in full:

\begin{minted}[fontsize=\footnotesize,breaklines]{text}
The field "availableQuantity" is unresolvable at the following path:
 query {
  products {
   availableQuantity <--
  }
 }
This is because:
 - The root type field "Query.products" is defined in the following subgraph: "catalog".
 - The field "Product.availableQuantity" is defined in the following subgraph: "mosaic".
 - The entity ancestor "Product" in subgraph "catalog" has no accessible target entities (resolvable @key directives) in the subgraphs where "Product.availableQuantity" is defined.
 - The type "Product" is not a descendant of any other entity ancestors that can provide a shared route to access "availableQuantity".
\end{minted}

\index{satisfiability!the error is a query}%
The error is a query. That is the part worth stopping on, and it is what makes
this the friendliest error in the whole composition catalogue. The composer does
not tell you that a rule was violated at a schema coordinate; it hands you a
document a client could have sent and points at the field it would have had to
fail on. \code{renderSelectionSet} in that same package draws the path from the
root and marks the offender with \code{<--}, and a path too deep to print gets
truncated with a count of what was left out.

The four reasons underneath read in order. Where the walk started. Where the
field lives. Why it could not get from the first to the second. And a last line
ruling out the possibility that some other entity on the path could have
provided a route.

The other three errors are the same shape for \code{price}, \code{reviews} and
\code{averageRating}. Two of those return composite types, so their rendered
paths open a selection set:

\begin{minted}[fontsize=\footnotesize]{text}
The field "price" is unresolvable at the following path:
 query {
  products {
   price { <--
     ...
    }
  }
 }
\end{minted}

\subsection*{One thing the error does not tell you}

\index{satisfiability!which route it names}%
Catalog has four root fields that return a \code{Product}:

\begin{minted}{text}
products   browseProducts   productById   productBySku
\end{minted}

All four are equally broken by this change, and the error names only
\code{Query.products}.

So the message reports \emph{a} route that fails, not every route that fails. I
have not gone into the walker to find out whether that is the first route it
tried, the shortest, or something else, and I am not going to guess in print.
What matters when you are reading one of these at half past five is that fixing
the named path is not evidence that you have fixed the graph. Recompose and see
whether it says something different.

\index{satisfiability!and the entity's own key}%
The fix here is obvious because the break was surgical, and that is the one
unrealistic thing about the experiment. Nobody types \code{resolvable: false} by
accident. What produces this error in a real graph is duller and further away
from the message: a subgraph that declares fields on somebody else's entity and
never declared the key, or declared it on a field it does not publish.

\index{federation!silent failures composition cannot see}%
Worth saying now, because a chapter about a checker owes it. The third way to
get an unreachable entity is to have the key, publish it, and have no working
reference resolver behind it, which is the failure
chapter~\ref{ch:the-first-cut-extracting-catalog} spent a section on. That one
composes. It composes because the composer reads two text files and a reference
resolver is not in either of them, which is a limit of the whole approach rather
than a gap in this tool. Section~\ref{sec:ch09-composition-as-a-gate} is about
what has to check the things composition structurally cannot.

That is the check with the reputation. The errors you will actually spend time
on are elsewhere, and they are stranger.

\index{satisfiability|)}%
